% Chalkline handout: Guided Notes (preset guided-notes v1).
% Document 01JQD00000000000000000002X, version 1.
% Compile with LuaLaTeX, TeX Live 2025 or later: latexmk -lualatex main.tex
\DocumentMetadata{lang=en, pdfversion=2.0, pdfstandard=ua-2, tagging=on, tagging-setup={math/setup=mathml-SE}}
\documentclass[11pt]{article}
\usepackage[letterpaper,margin=1in,headheight=15pt]{geometry}
\usepackage{amsmath}
\usepackage{fontspec}
\usepackage{unicode-math}
\directlua{luaotfload.add_fallback("clfallback", {"Latin Modern Math:mode=node;", "file:FreeSerif.otf:mode=node;"})}
\setmainfont{Latin Modern Roman}[RawFeature={fallback=clfallback}]
\setmathfont{Latin Modern Math}
\usepackage{xcolor}
\definecolor{clBlue}{HTML}{1F4E9C}
\definecolor{clRed}{HTML}{B3261E}
\definecolor{clGreen}{HTML}{1E6B34}
\definecolor{clOrange}{HTML}{A65200}
\definecolor{clPurple}{HTML}{6B3FA0}
\definecolor{clBrown}{HTML}{7A4A1E}
\usepackage{graphicx}
\usepackage{array}
\usepackage{booktabs}
\usepackage{tcolorbox}
\usepackage{fancyhdr}
\usepackage{lastpage}
\pagestyle{fancy}
\fancyhf{}
\renewcommand{\headrulewidth}{0.4pt}
\NewDocumentCommand{\clheader}{mm}{\fancyhead[L]{\small #1}\fancyhead[R]{\small #2}}
\NewDocumentCommand{\clfooter}{m}{\fancyfoot[C]{\small #1}}
\newif\ifcltitled
\NewDocumentCommand{\cltitle}{mm}{%
  \ifcltitled
    \begin{center}{\Large #1\par}\smallskip{#2\par}\end{center}%
  \else
    \title{#1}\author{#2}\date{}%
    \maketitle\thispagestyle{fancy}\vspace{-2em}\global\cltitledtrue
  \fi}
\NewDocumentCommand{\clquizheader}{mmm}{%
  \par\noindent #1: \rule{6cm}{0.4pt}\hfill #2: \rule{3cm}{0.4pt}\hfill #3: \rule{2cm}{0.4pt}\par\bigskip}
\NewDocumentEnvironment{cldefinition}{mm}{\begin{tcolorbox}[colback=white,colframe=black!55,boxrule=0.6pt,arc=2pt,left=6pt,right=6pt,top=4pt,bottom=4pt,before skip=8pt,after skip=8pt]\noindent\textbf{#1\IfBlankF{#2}{: #2}.}\ \ignorespaces}{\end{tcolorbox}}
\NewDocumentEnvironment{cltheorem}{mm}{\begin{tcolorbox}[colback=white,colframe=black!55,boxrule=0.6pt,arc=2pt,left=6pt,right=6pt,top=4pt,bottom=4pt,before skip=8pt,after skip=8pt]\noindent\textbf{#1\IfBlankF{#2}{: #2}.}\ \ignorespaces}{\end{tcolorbox}}
\NewDocumentEnvironment{clnote}{mm}{\par\medskip\noindent\textbf{#1\IfBlankF{#2}{: #2}.}\ \ignorespaces}{\par\medskip}
\NewDocumentEnvironment{clexample}{mm}{\par\medskip\noindent\textbf{#1\IfBlankF{#2}{: #2}.}\ \ignorespaces}{\par\medskip}
\NewDocumentEnvironment{clannouncement}{mm}{\par\medskip\noindent\textbf{#1\IfBlankF{#2}{: #2}.}\ \ignorespaces}{\par\medskip}
\NewDocumentEnvironment{clproblem}{mm}{\par\medskip\noindent\textbf{#1.}\IfBlankF{#2}{ \textit{(#2)}}\ \ignorespaces}{\par}
\NewDocumentEnvironment{clparts}{}{\begin{enumerate}[label=(\alph*)]}{\end{enumerate}}
\NewDocumentEnvironment{clsolution}{m}{\par\smallskip\noindent\textit{#1.}\ \ignorespaces}{\par\smallskip}
\NewDocumentCommand{\clanswer}{mm}{\par\smallskip\noindent\textit{#1:}\ #2\par}
\NewDocumentCommand{\clanswerboxed}{mm}{\par\smallskip\noindent\textit{#1:}\ \fbox{#2}\par}
\NewDocumentCommand{\clanswerspace}{m}{\par\vspace{#1\baselineskip}}
\NewDocumentCommand{\clplaceholder}{m}{\fbox{\textit{#1}}}
\NewDocumentCommand{\clblank}{m}{\par\noindent\fbox{\parbox[t][#1\baselineskip][t]{\dimexpr\linewidth-2\fboxsep-2\fboxrule\relax}{\mbox{}}}\par}
\NewDocumentCommand{\clcutline}{}{\par\vspace{1.5em}\noindent\makebox[\linewidth]{\dotfill}\par\vspace{1.5em}}
\NewDocumentEnvironment{clkeep}{}{\par\noindent\begin{minipage}{\linewidth}\setlength{\parskip}{0.5\baselineskip}}{\end{minipage}\par}
\NewDocumentEnvironment{clcolumns}{}{}{}
\setlength{\parindent}{0pt}
\setlength{\parskip}{0.5\baselineskip}
\usepackage{hyperref}
\hypersetup{hidelinks}
\hypersetup{pdftitle={Solving quadratic equations — Guided Notes}, pdfdisplaydoctitle=true}
\clheader{Algebra II — Period 3}{October 20, 2026}
\clfooter{Page \thepage\ of \pageref{LastPage}}

\begin{document}
\cltitle{Solving quadratic equations — Guided Notes}{Algebra II — Period 3 · José Nguyễn · October 20, 2026 · Unit 4: Quadratics}

\section*{Quadratic equations}

% blk:01JQD000000000000000000002
Today we solve equations of the form \(ax^{2} + bx + c = 0\) where \(a \neq 0\).

% blk:01JQD000000000000000000006
\begin{cldefinition}{Definition}{Quadratic equation}
% blk:01JQD000000000000000000004
A \textbf{quadratic equation} is a polynomial equation of degree 2.

\end{cldefinition}

% blk:01JQD00000000000000000000A
\begin{cltheorem}{Theorem}{Quadratic formula}
\clblank{2}
\end{cltheorem}

% blk:01JQD00000000000000000000M
\begin{clexample}{Example}{}
% blk:01JQD00000000000000000000C
Solve \(x^{2} - 5x + 6 = 0\).

% blk:01JQD00000000000000000000J
\begin{clsolution}{Solution}
% blk:01JQD00000000000000000000E
\begin{align*}
x^{2} - 5x + 6 &= 0 \\
(x - 2)(x - 3) &= 0
\end{align*}
% blk:01JQD00000000000000000000G
So \(x = 2\) or \(x = 3\).

\end{clsolution}
\end{clexample}

% blk:01JQD00000000000000000000R
\begin{clannouncement}{Announcement}{}
% blk:01JQD00000000000000000000P
Quiz Friday: sections 4.1–4.3. Bring a calculator \& a pencil!

\end{clannouncement}

% blk:01JQD00000000000000000000T
\begin{center}
\tagpdfsetup{table/header-rows={1}}
\begin{tabular}{ll}
\toprule
\(x\) & \(x^{2} - 5x + 6\) \\
\midrule
\(1\) & \(2\) \\
\(2\) & \(0\) \\
\(3\) & \(0\) \\
\bottomrule
\end{tabular}
\end{center}

% blk:01JQD00000000000000000000W
\begin{center}
\includegraphics[width=0.6\linewidth,alt={Parabola opening upward that crosses the x-axis at 2 and 3}]{assets/parabola-crop.png}
\end{center}

\section*{Practice}

% blk:01JQD00000000000000000001Q
\begin{clproblem}{1}{}
% blk:01JQD000000000000000000011
Solve by factoring.

\begin{clparts}
% blk:01JQD00000000000000000001B
\item
% blk:01JQD000000000000000000013
\[
x^{2} + 7x + 12 = 0
\]
% blk:01JQD000000000000000000017
\begin{clsolution}{Solution}
% blk:01JQD000000000000000000015
\[
(x + 3)(x + 4) = 0
\]
\end{clsolution}
% blk:01JQD000000000000000000019
\clanswerboxed{Answer}{\(x = -3\) or \(x = -4\)}
% blk:01JQD00000000000000000001N
\item
% blk:01JQD00000000000000000001D
\[
x^{2} - 9 = 0
\]
% blk:01JQD00000000000000000001H
\begin{clsolution}{Solution}
% blk:01JQD00000000000000000001F
\[
(x - 3)(x + 3) = 0
\]
\end{clsolution}
% blk:01JQD00000000000000000001K
\clanswerboxed{Answer}{\(x = \pm 3\)}
\end{clparts}
\end{clproblem}

% blk:01JQD000000000000000000021
\begin{clproblem}{2}{5 points}
% blk:01JQD00000000000000000001S
Use the quadratic formula to solve \(2x^{2} - 4x - 1 = 0\).

% blk:01JQD00000000000000000001X
\begin{clsolution}{Solution}
% blk:01JQD00000000000000000001V
\begin{align*}
x &= \frac{4 \pm \sqrt{16 + 8}}{4} \\
x &= \frac{2 \pm \sqrt{6}}{2}
\end{align*}
\end{clsolution}
% blk:01JQD00000000000000000001Z
\clanswerboxed{Answer}{\(x = 1 \pm \frac{\sqrt{6}}{2}\)}
\end{clproblem}

% blk:01JQD000000000000000000027
\begin{clproblem}{3}{}
% blk:01JQD000000000000000000023
\[
3x^{2} = 27
\]
% blk:01JQD000000000000000000025
\clanswerboxed{Answer}{\(x = \pm 3\)}
\end{clproblem}

% blk:01JQD00000000000000000002N
\begin{clproblem}{4}{}
% blk:01JQD000000000000000000029
A ball is thrown upward; its height is \(h(t) = -4.9t^{2} + 19.6t\) metres. When does it land?

% blk:01JQD00000000000000000002H
\begin{clsolution}{Solution}
% blk:01JQD00000000000000000002B
Set \(h(t) = 0\):

% blk:01JQD00000000000000000002D
\[
t(-4.9t + 19.6) = 0
\]
% blk:01JQD00000000000000000002F
so \(t = 0\) or \(t = 4\).

\end{clsolution}
% blk:01JQD00000000000000000002K
\clanswerboxed{Answer}{\(t = 4\) s}
\end{clproblem}

% blk:01JQD00000000000000000002Q
\begin{center}
\includegraphics[width=0.8\linewidth,alt={Handwritten work from the board, kept as an image}]{assets/crop-illegible.png}
\end{center}

% blk:01JQD00000000000000000002S
\begin{enumerate}
\item Factor when the roots are integers.
\item Otherwise use the formula \(x = \frac{-b \pm \sqrt{\Delta}}{2a}\).
\item Check by substitution.
\end{enumerate}

% blk:01JQD00000000000000000002V
Grading: 50\% homework \& 50\% quizzes; \#2 pencils; file\_name \textasciitilde{} \$5 fee \{braces\} \textbackslash{} back \textasciicircum{}\_\textasciicircum{}

\end{document}
